(a) By \exref{II.3.18}(a), a constructible subset is a finite union of locally closed subsets. Give each its reduced induced scheme structure. Its inclusion is of finite type, since the ambient scheme is noetherian, so it suffices to consider $f(X)$. A finite affine cover of $Y$ and finite affine covers of its inverse images reduce to affine $X,Y$. Replace $X$ by the reduced structures on its finitely many irreducible components. For each component, replace $Y$ by the reduced closure of its image; the morphism factors through this closure by \exref{II.3.11}(d), remains of finite type, and is dominant. Both schemes are now affine, integral, and noetherian.

(b) Prove the algebraic assertion by induction on the number of generators of $B$ over $A$. For $B=A$, take $a=b$. First consider $B=A[x]$ with one generator. If $x$ is transcendental over $\operatorname{Frac}A$, choose a nonzero coefficient $a$ of the polynomial $b$. Whenever $\varphi(a)\ne0$, the image of $b$ is a nonzero polynomial over $K$, so choose a value of $x$ at which it does not vanish.

If $x$ is algebraic over $\operatorname{Frac}A$, let $q$ be its monic minimal polynomial. In the field $(\operatorname{Frac}A)[x]$, express $b^{-1}$ as a polynomial in $x$. Choose $a\ne0$ clearing the coefficients of this expression and of $q$. Then $B_a=A_a[T]/(q)$ and $b$ is a unit in $B_a$. Every $\varphi$ with $\varphi(a)\ne0$ extends to $A_a$; choose a root of the monic polynomial $\varphi(q)$ in $K$ to extend it to $B_a$. The image of $b$ is nonzero because $b$ is a unit.

For the induction step, apply the one-generator case over $A[x_1,\ldots,x_{n-1}]$ to obtain a nonzero element $c$ of this intermediate ring. Apply the induction hypothesis to $c$, and then extend once more using the one-generator case. This proves the assertion.

Apply it with $b=1$. For every $\mathfrak p\in D(a)\subseteq\operatorname{Spec}A$, embed $k(\mathfrak p)$ into an algebraically closed field and compose with $A\to k(\mathfrak p)$. The extension to $B$ has prime kernel contracting to $\mathfrak p$. Thus $D(a)\subseteq f(X)$.

(c) Use noetherian induction on closed subsets of $Y$. For a morphism into $Y$, if the closure of its image is proper, apply induction there. Otherwise some irreducible component of $X$ has dense image in an irreducible component of $Y$; the reductions of (a) and (b) show, on an affine open of that component of $Y$, that the image contains a nonempty open subset. If $Y$ is irreducible this is a nonempty open $U\subseteq Y$, and $f(X)\setminus U$ is the image of $f^{-1}(Y\setminus U)\to Y\setminus U$, constructible by induction. Hence $f(X)$ is constructible. If $Y$ is reducible, apply induction to its proper irreducible components and the corresponding inverse images, and take their finite union. Together with (a), this proves the theorem.

(d) The morphism $\mathbb A_k^2\to\mathbb A_k^2$, $(s,t)\mapsto(s,st)$, has image $D(x)\cup\{(0,0)\}$. This image is dense and proper, so is not closed. It is not open, since its intersection with the line $x=0$ is the single point $(0,0)$.
